\subsection{European Natural Gas Markets}\label{sec:gas-markets}
Natural gas is a physical commodity that exists within a gas grid. The TTF is a virtual trading hub in the Netherlands. A virtual trading hub allows participants to transfer ownership without reference to a specific physical location. TTF is operated by Gasunie Transport Services (GTS) within the Dutch gas transmission network. Gas futures contracts referencing TTF are traded on ICE Endex, the European energy exchange operated by Intercontinental Exchange. Futures contracts are described in Sec.~\ref{sec:futures}.

TTF is the most liquid natural gas market in Europe, with traded volumes 4.5 times the combined volume of the other eight main European hubs \parencite{heatherEuropeanTradedGas2026}. Therefore, the TTF gas price serves as the reference price for gas contracts. 

At virtual hubs such as TTF, gas price formation is determined by short-run supply and demand \parencite{obadiDRIVINGFUNDAMENTALSNATURAL2020}. In the short run the infrastructure of the gas grid is fixed, and the supply of gas cannot expand in response to a change in demand. Therefore, the adjustment takes place through the price, rather than the quantity.

Gas demand originates in three sectors. The first source is residential and commercial heating, which makes consumption seasonal.
Deviations of weather from the seasonal norm also affect prices \parencite{brownWhatDrivesNatural2008}. Industrial production is the second source of demand. In the chemical sector, for example, gas is used as an energy source and as a raw material. The third source is power generation, where gas competes with coal. As a result, the gas market is linked to the coal market and, through emission allowances, to the carbon market.

On the supply side, natural gas prices depend on pipeline flows and deliveries of liquefied natural gas (LNG). LNG is natural gas cooled to liquid form for transport by ship. 

As has been observed repeatedly in the past, the European natural gas market is highly sensitive to supply disruptions. The European gas price crisis between 2021 and 2022 is the most severe episode. During the COVID-19 pandemic in 2020, energy demand collapsed. This was followed by economies reopening, and gas demand recovered faster than production could follow. From mid-2021, the post-pandemic increase in demand was also combined with reduced Russian pipeline flows, strong Asian competition for LNG cargoes and storage levels below the seasonal norm \parencite{heather2022unfortunate}. As a result, TTF gas prices rose from roughly 18 EUR/MWh in the first quarter of 2021 to a peak of 339 EUR/MWh on 26 August 2022 and then declined from early 2023. 

The crisis eventually raised the question of whether the TTF price benchmark had become less representative of the underlying gas market. In November 2022, the European Commission proposed a Market Correction Mechanism (MCM), which would cap gas prices during crisis periods \parencite{MCM-euro-commision}. The mechanism applied until January 2025 but was never activated, as gas prices declined after 2023.

A more recent disruption followed the military campaign against Iran by Israel and the United States on 28 February 2026. Iran responded with attacks on shipping and energy infrastructure in the Gulf. Between 27 February and 2 March, the traffic through the Strait of Hormuz fell from 95 vessels to 4 vessels per day \parencite{persian-gulf-26}. These events resulted in a spike of TTF gas prices from 27 EUR/MWh on 5 January 2026 to 61 EUR/MWh on 19 March 2026.

\subsection{Futures Contracts}\label{sec:futures}
A derivative is a contract whose value depends on the price of another asset, called the underlying. A forward contract is an agreement between two parties to buy or sell an asset at a fixed price on a specified future date. A futures contract is a standardized forward contract traded on an exchange \parencite[28--30]{hullOptionsFuturesOther2022}.
The quantity and the delivery date are set by the exchange. The price is agreed between the parties and remains fixed for the contract. As a result, all contracts for a given month are identical in their terms and can be freely transferred. TTF gas futures are quoted in EUR/MWh and are physically delivered, with delivery spread evenly across the hours of the delivery month. Futures contracts are settled daily against a settlement price published by the exchange. The settlement price is determined by the exchange at the end of each trading day and is not necessarily a price at which a transaction occurred. The daily gain or loss is exchanged in cash between the two parties. 

The payoff of a long position in the futures contract is the difference between the settlement price $F_T$ at expiry $T$ and the agreed price $K$. The payoff of a short position is the difference between the agreed price $K$ and the settlement price $F_T$. Both payoffs can be positive or negative \parencite[29]{hullOptionsFuturesOther2022}.

Price formation of natural gas in Europe has generally taken place on the futures rather than spot market \parencite{obadiDRIVINGFUNDAMENTALSNATURAL2020}. In a spot market the gas is traded for immediate delivery. A futures contract can be traded repeatedly before expiry. As a result, the market is open for hedgers and speculators without an interest in the physical commodity.

At any point in time, contracts for multiple consecutive months are listed and can be traded in parallel. The contract with the nearest delivery on a given day is the front-month contract (M1), and the next is the second-month contract (M2). Trading concentrates in the front-month contract, hence it is commonly used as the reference for the current gas price. Since each contract ceases trading before its delivery month, the contract that is the front month changes over time. The prices of the contracts sorted by the delivery month form the term structure of the gas market. The market is in contango if the term structure is upward sloping, and in backwardation if it is downward sloping. The price difference between two consecutive months is the calendar spread. A calendar spread describes the shape of the term structure. 